% =============================================================================
% INTRODUCTION TO THE EU READING SUMMARY TEMPLATE
% =============================================================================
% Copy this file to create new reading summaries. All summaries use this structure.
% Compile: pdflatex your-filename.tex
%
% Sections (in order): Rhetorical Framework | Bibliographic Info | Summary |
% Main Arguments | Context & Analysis | Relevance to Course | Discussion Questions
% =============================================================================

\documentclass[11pt,a4paper]{article}
\usepackage[utf8]{inputenc}
\usepackage[T1]{fontenc}
\usepackage[margin=2.5cm]{geometry}
\usepackage{booktabs}
\usepackage{enumitem}
\setlist{nosep,leftmargin=*}
\usepackage{parskip}
\usepackage{microtype}
\usepackage{hyperref}
\usepackage{xcolor}
\definecolor{eublue}{RGB}{0,51,153}

\hyphenation{Spitzenkandidat Spitzenkandidaten}
\hypersetup{
  colorlinks=true,
  linkcolor=eublue,
  urlcolor=eublue,
  pdftitle={Reading Summary},
  pdfauthor={Introduction to the EU, UCM}
}

\begin{document}

% --- TITLE: Author: Title (italic) \ Subtitle
\title{Author Name: \textit{Work Title}\\Key Themes \& Arguments --- Rhetorical Analysis}
\author{Introduction to the European Union\\Universidad Complutense de Madrid\\Week N}
\date{}
\maketitle
\vspace{-1em}
\hrule
\vspace{1em}

% --- 1. RHETORICAL FRAMEWORK
\section*{Rhetorical Framework (Ethos, Logos, Pathos)}

\begin{table}[h]
\centering
\begin{tabular}{@{}lll@{}}
\toprule
\textbf{Appeal} & \textbf{Meaning} & \textbf{Author's Use} \\
\midrule
\textbf{Ethos} & Credibility, authority, trust & \\
\textbf{Logos} & Logic, structure, argument & \\
\textbf{Pathos} & Emotion, stakes, persuasion & \\
\bottomrule
\end{tabular}
\end{table}

\textbf{Key insight:} One paragraph on how the author blends ethos, logos, pathos.

% --- 2. BIBLIOGRAPHIC INFORMATION
\section*{Bibliographic Information}

\begin{itemize}
\item \textbf{Author:}
\item \textbf{Title:}
\item \textbf{Source:}
\item \textbf{Year:}
\end{itemize}

% --- 3. SUMMARY & KEY POINTS
\section*{Summary \& Key Points}

\subsection*{I. First Theme}
Content\ldots

\subsection*{II. Second Theme}
Content\ldots

% --- 4. MAIN ARGUMENTS
\section*{Main Arguments}

\subsection*{I. First Argument}
Content\ldots

\subsection*{II. Second Argument}
Content\ldots

% --- 5. CONTEXT & ANALYSIS
\section*{Context \& Analysis}

\begin{itemize}
\item \textbf{Historical:}
\item \textbf{Connections:}
\end{itemize}

% --- 6. RELEVANCE TO COURSE
\section*{Relevance to Course}

\begin{itemize}
\item \textbf{Ethos:}
\item \textbf{Logos:}
\item \textbf{Pathos:}
\item \textbf{All three:}
\end{itemize}

% --- 7. DISCUSSION QUESTIONS
\section*{Discussion Questions}

\begin{enumerate}
\item Question 1?
\item Question 2?
\item Question 3?
\end{enumerate}

% --- FOOTER
\vspace{1em}
\noindent\footnotesize\textit{Introduction to the European Union, Universidad Complutense de Madrid.}

\end{document}
